% Lösungen Einheit 4 von 4 – Lot und Extremum (zusammenbau v0.1)
\einheitenkopf{Einheit 4 von 4 · Lot und Extremum}

\erg{18}{a) $\overrightarrow{PF}$ und der Richtungsvektor von $g$ – das Lot steht senkrecht auf der Geraden \quad b) Kürzeste Verbindung ist das Lot: $\overrightarrow{PX_\lambda} \circ (2 | 1 | 2) = 0$ mit $X_\lambda$ auf $g$; das liefert $9 - 9\lambda = 0$, $\lambda = 1$, $F(2 | 2 | 2)$, Abstand $|\overrightarrow{PF}| = \sqrt{5} \approx 2{,}24$ \quad c) Der Spitzenwinkel ist am größten, wenn die Höhe $|MF_k|$ von der Basismitte $M(0 | 0 | 2)$ am kleinsten ist; die kürzeste Verbindung ist das Lot auf die Gerade der Punkte $F_k$ (Richtungsvektor $(0 | 1 | -2)$): $\overrightarrow{MF_k} \circ (0 | 1 | -2) = (0 | k | 10 - 2k) \circ (0 | 1 | -2) = k - 20 + 4k = 5k - 20 = 0$, also $k = 4$ \quad d) Die Dreiecke $ABM$ (rechter Winkel bei $B$) und $AFB$ (rechter Winkel bei $F$) haben den Winkel bei $A$ gemeinsam, sind also ähnlich. In $ABM$ ist $|AB| : |BM| = 2 : 1$, weil $M$ die Seite halbiert; dasselbe Verhältnis gilt für die entsprechenden Seiten $|AF| : |BF|$ in $AFB$.}
\erg{19}{a) Mit $M(0 | 3 | 3)$ als Mittelpunkt von $RS$ ist der Flächeninhalt $\frac{1}{2} \cdot |RS| \cdot |MT_\lambda|$; $|RS|$ ist fest, also ist die Fläche am kleinsten, wenn $|MT_\lambda|$ am kleinsten ist – das ist das Lot von $M$ auf $t$, und dort steht $\overrightarrow{MT_\lambda} = (1 + 2\lambda | 0 | \lambda - 3)$ senkrecht auf dem Richtungsvektor $(2 | 0 | 1)$; Lösung $5\lambda - 1 = 0$, $\lambda = 0{,}2$}
\erg{20}{a) Fehler: Die Fläche hängt bei fester Basis von der Höhe ab, nicht vom Abstand der Spitze zu einem Endpunkt. Richtig: $|MT_\lambda|$ mit dem Mittelpunkt $M$ von $RS$ minimieren – die Lotbedingung, $\overrightarrow{MT_\lambda}$ senkrecht zum Richtungsvektor von $t$ \quad b) Jede andere Strecke von $P$ zu einem Punkt $X$ auf $g$ ist Hypotenuse im rechtwinkligen Dreieck $PFX$ mit dem Lotfußpunkt $F$; die Hypotenuse ist länger als die Kathete $PF$. \quad c) Kürzester Weg ist das Lot: $(2 - 2\lambda | 4 - \lambda | 0) \circ (2 | 1 | 0) = 4 - 4\lambda + 4 - \lambda = 8 - 5\lambda = 0$, $\lambda = 1{,}6$: Auffahrt bei $(3{,}2 | 2{,}6 | 0)$, Länge $\sqrt{1{,}44 + 5{,}76} = \sqrt{7{,}2} \approx 2{,}68$, also $\approx 268$ m}
